\section{域外司法协助}
\begin{enumerate}
  \item 文书送达。9 种途径。一出一进。
        \begin{kk}{出}
          \zx{条约外交使领馆，诉代分办后要权。传真电邮能收悉，董监高公告兜底。进基本同出}。
          \begin{enumerate}
            \item 国际条约。海牙公约。\zx{使领馆} \yjt \bs{司法部}！！ \yjt \zx{最高法}。
            \item 外交途径。
                  \begin{enumerate}
                    \item \bs{没有条约关系才可以，不能既条约又外交途径}。
                    \item 无需外交部同意。\al{只有经外交部同意，甲国才能向其公民送达传票 \cuo}。
                  \end{enumerate}
            \item 使领馆。只对\zd{其本国公民}。\sjt 只能域外。\xjt 称\zd{通用送达}途径。\al{通过我国使领馆向甲国乙公民送达 \cuo}。
            \item 诉讼代理人、任何代表机构。\zd{“代字”}。\al{应当向代表处送达文书 \cuo{是可以}}。
            \item \zd{被授权的}分支机构、业务代办人。
            \item \zd{邮寄且允许邮寄}。超 3 月失败。
            \item 传真、电邮要\zd{确认能收悉}。
            \item 在我国的法定代表人、主要负责人（董监高）。
            \item 公告兜底途径。3 月成功。
            \item 中国驻美国使领馆送达文书给中国当事人，现在需要\zd{经过所在国允许 \hz 所在国同意}。新增。
          \end{enumerate}
          \tcblower
          \bs{排除公告}，其他\zd{无顺序}要求、\zd{可同时}多种、按\zd{最先送达}的认定日期。
        \end{kk}

        \begin{kk}{进}
          \begin{enumerate}
            \item 条约 \dy 外交途径。
            \item 外国驻华使领馆可以向该国公民送达。
          \end{enumerate}
        \end{kk}

        \begin{kk}{海牙送达公约}
          \begin{enumerate}
            \item 外国法院 \yjt \zx{该国驻华使领馆} \yjt \bs{司法部}！！ \yjt \zx{最高法} \yjt 有关法院 \yjt 受送达人。
            \item 以下理由，被请求方法院\zx{不得拒绝送达}。
                  \begin{enumerate}
                    \item \zd{有关期限已过}。
                    \item 为中国法院\zd{专属管辖} \hz 不承认对该事项提起诉讼的权利。
                    \item 未附有被请求方文字译本，但附有英文 \hz 法文译本。
                  \end{enumerate}
            \item 可以拒绝的。\bs{不属于职权范围的、危害国家安全}。
            \item 受送达人\zd{有权拒收未附有中文译本}的司法文书。
          \end{enumerate}
        \end{kk}
  \item 域外调取证据。我国只有 2 种。
        \begin{enumerate}
          \item 代为取证。以\zx{条约为基础} + \bs{请求主体是司法机关} + \zx{取证请求书}。\al{取证委托书 \cuo}。
          \item 使领馆取证。同使领馆送达，\zd{只能对其本国公民}且不得采取强制措施。
          \item 我国原则上\zx{不允许特派员取证、自行取证}。自行 = 自己跑过去、自己去申请、自己对接。
          \item 海牙取证公约。
                \begin{enumerate}
                  \item 不允许直接传递，以\bs{司法部}！！中介。
                  \item 因为是司法协助，所以\bs{仅限司法程序证据}，禁止调取行政程序的证据。
                \end{enumerate}
        \end{enumerate}
  \item 外国法院判决的承认和执行。当事人 \hz \zd{外国法院}申请，其他基本同仲裁。
        \begin{kk}{要求}
          \zx{协议互惠已生效，公序不专不缺席}。
          \begin{enumerate}
            \item 有双边司法协助\zd{协议互惠}关系。\bs{例外离婚判决，但也仅限离婚}。\al{执行解除夫妻身份关系之外的 \cuo}。
            \item 判决\zd{已生效}、不违反\zd{公序}良俗、\zd{不}属于我国法院\zd{专}属管辖。
            \item 原则上拒绝承认缺席判决 \hz 裁定。例外。\zx{已合法传唤 + 判决 \hz 裁定已对此说明}。
          \end{enumerate}
        \end{kk}
  \item 仲裁的承认和执行。1958 年承认及执行外国仲裁裁决公约/纽约公约。
        \begin{enumerate}
          \item \zx{我国做了 4 个保留}。
                \begin{enumerate}
                  \item 仲裁地在缔约国领土。
                  \item 仅限商事仲裁（合同、侵权、财产权益纠纷），\zd{不包括民事仲裁裁决}。
                  \item 临时仲裁裁决也可进行申请承认和执行。
                  \item 得是私法主体之间，任何一方\bs{不能是主权国家}。
                \end{enumerate}
          \item \bs{仅限当事人}申请。
          \item 受理法院。\zx{被执行人住所地和被执行财产所在地的中院}。同民诉。
          \item 裁决生效后 2 年内提。\zd{被驳回可起诉}。救济和港澳台有细微差别。
                \begin{enumerate}
                  \item 可同时、可分开提。
                  \item 申请承认的，\zd{合议庭}仅裁定是否承认 \yjt 承认 \yjt 申请执行的 2 年期间自裁定生效之日起\zd{中断}。
                  \item 承认、执行\bs{裁定}中院\zd{可以直接作出}。\zx{不承认需要履行内部报核程序}。\al{直接裁定拒绝承认、执行 \cuo}。
                \end{enumerate}
          \item 限制性规定。
                \begin{enumerate}
                  \item 对外国仲裁裁决\zx{无权撤销，只能作拒绝承认 \hz 不予执行裁定} \yjt 法院内部报核。
                  \item 原则上\zx{无权主动审查，只审程序}。例外。公序良俗、是我国加入时保留的事项。
                \end{enumerate}
        \end{enumerate}
\end{enumerate}
